\section{Conclusion}
\label{sec:conclusion}

\tothink{full conclusion prose; outline below.}

\noindent We presented an adaptive conservative selection policy for hybrid
retrieval over the \blackrim multi-agent memory pool. The policy treats
fusion-config selection as a contextual bandit with discrete arms,
query class as context, per-query NDCG@10 as reward, and a per-class
baseline (BM25 / BM25+decay) as the conservative-bandit safety
constraint. The architecture stacks BM25 + SPLADE + dense + RRF +
decay + MMR + optional cross-encoder reranker, with a SPLADE
corpus-size gate and a Critical-tier reranker hook. Empirical
evaluation \tothink{recap headline numbers once data accumulates}
shows the policy matches the per-class best of any baseline without
regressing below the per-class baseline floor.

\noindent The paper's central methodological claim is that
\emph{retrieval-policy selection in agent-memory systems is a conservative
contextual-bandit problem with per-class baselines}, structurally
distinct from the tier-cost-ratio framing of LLM model routing. The
per-class no-regression constraint is the operational lever; classical
IR evaluation (NDCG@10, MRR@10, recall@k, paired bootstrap) provides
the measurement infrastructure; the conservative-bandit theory
\cite{wu2016conservative,kazerouni2017conservative} provides the
selection rule.

\paragraph{Acknowledgements.}
\tothink{contributors, operators, and the open-source artefacts. Same
shape as MOA-1c.}
